\renewcommand{\arraystretch}{1.3}
\captionsetup[table]{justification=raggedright,singlelinecheck=false}
\begin{table}[H]
\caption{Rancangan Tiga Lapis Validasi YAML}
\label{tbl:yaml_validation_iv}
\centering
\begin{tabular}{|p{0.08\textwidth}|p{0.20\textwidth}|p{0.25\textwidth}|p{0.35\textwidth}|}
\hline
\textbf{Lapis} & \textbf{Alat} & \textbf{Jenis Pemeriksaan} & \textbf{Aksi jika Gagal} \\
\hline
1 & \texttt{yaml.safe\_load} & Sintaksis YAML (format, indentasi, tipe data) & Hentikan validasi; kembalikan \texttt{syntax\_errors} \\
\hline
2 & \texttt{kubernetes- validate} v1.29 & Kepatuhan skema (nama \textit{field}, tipe, struktur) & Tambahkan \texttt{schema\_errors}; lanjut ke Lapis 3 \\
\hline
3 & Kueri Cypher Neo4j & Ketersediaan \textit{required fields} (\texttt{is\_required: true}) & Tambahkan \texttt{missing\_fields}; tandai YAML tidak valid \\
\hline
\end{tabular}
\end{table}
